\section{Einleitung}\label{sec:einleitung}

\subsection{Motivation und Projektkontext}

Diese Hausarbeit entstand als Ergänzung zur Bachelorarbeit über Counterfactual Regret Minimization (CFR) zur Lösung von Poker-Spielen. Während die Bachelorarbeit sich auf die Entwicklung und das Training von KI-Strategien konzentrierte, wurde hier eine grafische Benutzeroberfläche entwickelt, die die bereits implementierte Spiel-Logik aus der Bachelorarbeit nutzt. Das Ziel war es, die trainierten Strategien für menschliche Spieler zugänglich zu machen und gleichzeitig die Möglichkeit zu schaffen, dass zwei menschliche Spieler über ein Netzwerk gegeneinander spielen können.

Das Projekt umfasst zwei verschiedene Spielmodi, die unterschiedliche Anforderungen an die Architektur stellen. Der erste Modus ist \textbf{Agent vs Human}, bei dem der Spieler gegen trainierte KI-Strategien antritt, die in der Bachelorarbeit entwickelt wurden. Diese Strategien basieren auf CFR-Algorithmen und wurden für verschiedene Poker-Varianten wie Kuhn Poker und Leduc Hold'em trainiert. Die Herausforderung hierbei liegt in der Integration der trainierten Strategien, die als komprimierte Pickle-Dateien vorliegen, und der automatischen Spielsteuerung des Agenten. Der zweite Modus ist \textbf{Human vs Human}, ein Multiplayer-Spiel über Netzwerk, bei dem zwei menschliche Spieler über einen WebSocket-Server gegeneinander spielen können. Diese Variante erfordert eine Client-Server-Architektur mit Echtzeit-Kommunikation und Thread-Safety, da mehrere Clients gleichzeitig auf die Spiel-Logik zugreifen.

Die Implementierung demonstriert verschiedene fortgeschrittene Python-Konzepte, die über die Grundlagen hinausgehen. PyQt6 wird für die GUI-Entwicklung verwendet und zeigt das Signal-Slot-Pattern als Implementierung des Observer-Patterns. Asyncio wird für asynchrone Netzwerk-Kommunikation eingesetzt, wobei die Integration mit der Qt Event Loop eine besondere Herausforderung darstellt. Threading wird für Thread-Safety bei Multiplayer-Szenarien verwendet, pytest ermöglicht umfassende Tests ohne GUI-Fenster, und Design Patterns wie das Observer-Pattern durch Signals und Slots werden praktisch angewendet.

\subsection{Zielsetzung}

Die Hauptziele dieser Arbeit waren die Entwicklung einer funktionsfähigen Poker-GUI mit zwei Spielmodi, die Wiederverwendung der Game-Logik aus der Bachelorarbeit ohne Änderungen an der bestehenden Implementierung, die Integration von trainierten KI-Strategien für den Agent vs Human Modus, die Implementierung eines WebSocket-Servers für Multiplayer-Echtzeit-Kommunikation, umfassende Tests mit pytest, die ohne GUI-Fenster ausgeführt werden können, sowie die Demonstration von Advanced Python Konzepten wie Signals/Slots (Observer-Pattern), asyncio+QThread Integration, Thread-Safety mit \texttt{threading.Lock}, Decorators für Rate Limiting und Pickle sowie Kompression für Strategie-Serialisierung.

Ein wichtiger Aspekt war dabei, dass beide Spielmodi die gleichen UI-Komponenten verwenden sollten, um Code-Duplikation zu vermeiden und die Wartbarkeit zu erhöhen. Zudem sollte die Architektur so gestaltet sein, dass die Spiel-Logik aus der Bachelorarbeit ohne Änderungen wiederverwendet werden kann, was eine klare Trennung zwischen GUI-Schicht, Client-Schicht, Server-Schicht und Logik-Schicht erforderte.

\subsection{Überblick über die Arbeit}

Diese Arbeit gliedert sich in sieben Kapitel. Kapitel 2 stellt den Hintergrund und theoretischen Rahmen vor und diskutiert relevante Technologien und Frameworks sowie ähnliche Ansätze. Dabei werden PyQt6, WebSocket, asyncio, pytest und andere verwendete Technologien vorgestellt und begründet. Kapitel 3 beschreibt das Konzept und die Umsetzung mit Fokus auf Advanced Python Konzepte, erläutert die Architektur, Designentscheidungen, Code-Struktur sowie GUI-, Server- und Client-Komponenten im Detail. Kapitel 4 enthält die Developer-Dokumentation mit Installationsanleitung, Systemvoraussetzungen und benötigten Bibliotheken. Kapitel 5 stellt die Test-Strategie und die implementierten Tests vor, die das Verhalten des HTTP-Servers, WebSocket-Servers und der Game-Logik validieren. Kapitel 6 präsentiert die Ergebnisse und Evaluation, diskutiert was funktioniert und was nicht, sowie Messungen und Beobachtungen. Kapitel 7 fasst die Diskussion und das Fazit zusammen, bewertet die technische Umsetzung, diskutiert die erlernten Advanced Python Konzepte und gibt einen Ausblick auf mögliche Erweiterungen. Code-Beispiele zu wichtigen Konzepten finden sich im Anhang.
